\section{线性回归从投影到概率模型}

假设我们想根据房屋面积、楼龄和地段估计租金。最先想到的模型可能是

\[
\widehat y=\beta_0+\beta_1x_1+\beta_2x_2+\beta_3x_3.
\]

它的价值不在于断言世界真是一张平面，而在于给出一个可检验的近似：保持其他变量不变时，每个特征对预测值贡献多少；模型剩下的误差有多大；加入更复杂结构前，最朴素的规律已经能解释多少。

\subsection{为什么``画一条最贴近数据的线''还不够}

若只看二维散点图，``贴近''似乎很直观。进入高维以后，必须先说明怎样衡量误差。最小二乘选择平方残差和：

\[
\widehat\beta\in\argmin_\beta\|y-X\beta\|_2^2,
\]

其中 \(X\in\mathbb R^{n\times p}\) 的每一行是一条样本，每一列是一项特征，\(y\) 是响应向量。截距可以通过给 \(X\) 增加一列全 1 向量表示。

平方损失不是唯一可能，却有一个极清楚的几何解释：所有可能的拟合值 \(X\beta\) 都位于 \(X\) 的列空间中，最小二乘是在这个空间里找离 \(y\) 最近的点。最优残差

\[
r=y-X\widehat\beta
\]

必须与列空间正交，因此

\[
X^\top r=0,\qquad
X^\top X\widehat\beta=X^\top y.
\]

后一个式子叫正规方程。它说明了答案满足什么关系，却不意味着计算时应显式求 \((X^\top X)^{-1}\)。形成 \(X^\top X\) 会把条件数大致平方；实际求解通常优先使用 QR 分解，近秩亏时则使用 SVD。

\subsection{一个能看见投影的例子}

把三个样本的特征中心化为 \(x=(-1,0,1)^\top\)，响应为 \(y=(1,2,2)^\top\)。含截距的一元回归给出

\[
\widehat\beta_0=\bar y=\frac53,\qquad
\widehat\beta_1=\frac{x^\top y}{x^\top x}=\frac12.
\]

于是拟合值为

\[
\widehat y=\left(\frac76,\frac53,\frac{13}6\right)^\top,
\]

残差为 \((-1/6,1/3,-1/6)^\top\)。残差之和为零，也与 \(x\) 的内积为零。这两项检查分别对应``与截距列正交''和``与特征列正交''。它们不是偶然的数值现象，而是一阶最优条件在具体数据上的表现。

这个例子还暴露出一个容易忽略的区别：最优拟合值是唯一投影，但系数未必唯一。若两列特征完全相同，许多组系数都能产生同一个 \(X\beta\)。此时预测面仍可确定，单个系数却无法由数据区分。参数解释比预测需要更强的可识别性。

\subsection{从几何问题到概率模型}

仅为了求一个最小二乘拟合，不需要假设误差服从正态分布。若进一步写成

\[
Y=X\beta+\varepsilon,\qquad
\E(\varepsilon\mid X)=0,
\]

便开始讨论重复抽样下的估计性质。若再假设条件同方差且误差不相关，

\[
\operatorname{Var}(\varepsilon\mid X)=\sigma^2I,
\]

就能得到经典标准误。若还假设 \(\varepsilon\mid X\sim\mathcal N(0,\sigma^2I)\)，最小二乘又等价于最大似然估计，并可在有限样本下推导精确的区间与检验。

这些假设承担的任务不同：

\begin{itemize}
\tightlist
\item
  条件均值为零使 \(\beta\) 描述条件均值，并支撑无偏性；
\item
  方差结构决定标准误怎样计算；
\item
  正态性主要支撑小样本精确分布，而不是让最小二乘解``存在''；
\item
  因果解释还需要混杂控制、干预或识别假设，单靠回归公式不会自动产生。
\end{itemize}

因此``我成功拟合了一条线''\,``这个系数有统计不确定性保证''和``改变该变量会导致响应改变''是三种完全不同的陈述。

\subsection{一次完整建模不止求出系数}

设我们仍以租金预测为例。合理流程通常是：

首先明确部署时可获得哪些特征，并在切分训练集之前固定预测目标。类别编码、缺失值填补、标准化等操作只能在训练数据上拟合，否则测试信息会泄漏进模型。随后用 QR 或 SVD 求解，并同时观察：

\begin{itemize}
\tightlist
\item
  残差是否随拟合值呈漏斗形或曲线形；
\item
  少数高杠杆点是否主导了系数；
\item
  特征矩阵的条件数是否过大；
\item
  训练误差和独立测试误差是否差异悬殊；
\item
  部署人群是否超出了训练数据的特征范围。
\end{itemize}

如果残差出现稳定曲线，问题不是``优化器没收敛''，而是线性表示遗漏了结构；如果训练表现很好而测试表现差，可能是特征过多或切分方式错误；如果测试集来自同一时期，真实部署却发生在半年后，随机切分又可能高估泛化能力。

还应根据任务选择指标。均方误差会重罚大误差，绝对误差更抗离群值；若高价房预测错一万元与低价房错一万元的业务后果不同，甚至需要直接建模不对称成本。目标函数方便，并不等于目标函数代表真正关心的后果。

\subsection{边界：线性究竟指什么}

``线性模型''通常指对参数线性，而不要求对原始变量只能画直线。例如

\[
\widehat y=\beta_0+\beta_1x+\beta_2x^2
\]

仍是线性回归，因为待估参数以线性方式出现。加入多项式、样条或交互项可以扩展表示，但每加入一种自由度，都在增加方差和外推风险。模型是否合适要靠问题结构与独立数据判断，而不是靠``线性''这个名字。

\subsubsection{留给读者的半开放检验}

构造两列几乎相同但不完全相同的特征，分别用正规方程、QR 和 SVD 求最小二乘系数。逐渐缩小两列的差异，记录系数、拟合值和测试预测的变化。先预测哪一项最早变得不稳定，再用实验解释为什么。下一节将从这个失败出发引入正则化。

相关基础可回看``正交与最小二乘''、``数值线性代数''与``点估计''。
